% PM証明体系：公理・定理の証明まとめなおし
% 第8章（PM証明体系）は公理と原始規則を示すのみで，定理2以降の完全な証明が
% ほとんど無いため，A1〜A4と原始規則PDR1・PDR2だけから，定理1〜定理6a
% （対策プリントの問5で扱った7つ）とそこで使う派生規則をすべて導出し直す。
\documentclass[autodetect-engine,dvi=dvipdfmx,ja=standard]{bxjsarticle}

\makeatletter
\providecommand*{\input@path}{}
\g@addto@macro\input@path{{./}{../includes/}}
\makeatother

\input{protocol.tex}
\input{preamble.tex}

\title{PM証明体系：公理・定理の証明まとめ}
\author{論理学 第8章 対応}
\date{\today}

\begin{document}
\maketitle

第8章「PM証明体系」は公理 A1--A4 と原始導出規則 PDR1・PDR2 を提示するが，
そこから実際に定理を導く過程は定理2（反射律）にしか示されていない。
本資料は，\textbf{A1--A4 と PDR1・PDR2 だけを出発点として}，対策プリントの問5で
扱った7つの定理（PM-Th.1〜PM-Th.6a）と，途中で必要になる派生規則
（PDR3，DR4）を，省略なくすべて導出し直したものである。
「④SPC恒真式だからDR-E8で証明できる」という意味論的な近道は使わず，
すべて構文論的な証明（公理と規則の適用のみ）で統一する。

\section{出発点：公理・原始規則・定義}

\subsection*{公理}

\[
  \mathrm{A1}\qquad p\lor p\supset p
  \qquad\qquad
  \mathrm{A2}\qquad p\supset p\lor q
\]
\[
  \mathrm{A3}\qquad p\lor q\supset q\lor p
  \qquad\qquad
  \mathrm{A4}\qquad (p\supset q)\supset\{(r\lor p)\supset(r\lor q)\}
\]

付値\(V_i\)により，4つとも恒真式であることが確かめられる：
\(V_i[p]\)が\(1\)でも\(0\)でも\(V_i[p\lor p\supset p]=1\)（A1），
\(V_i[p]=1\)なら\(V_i[p\supset p\lor q]=1\)は自明で，\(V_i[p]=0\)なら前件が偽で自動的に真（A2），
\(\lor\)の対称性から\(V_i[p\lor q]=V_i[q\lor p]\)なので常に真（A3），
\(p\supset q\)が偽なら前件が偽で自動的に真，真なら\(V_i[p]=1\Rightarrow V_i[q]=1\)なので
\(r\lor p\)が真になるすべての場合に\(r\lor q\)も真になる（A4）。

\subsection*{原始導出規則}

\[
  \frac{\mathrm{PM}\vdash\alpha[p]}{\mathrm{PM}\vdash\alpha[\beta]}\quad\mathrm{PDR1（一様代入）}
  \qquad\qquad
  \frac{\mathrm{PM}\vdash\alpha\quad\mathrm{PM}\vdash\alpha\supset\beta}{\mathrm{PM}\vdash\beta}\quad\mathrm{PDR2（分離）}
\]

\subsection*{定義}

\[
  \alpha\supset\beta\;\mathrel{=_{\mathrm{df}}}\;\sim\alpha\lor\beta
  \qquad\qquad
  \alpha\land\beta\;\mathrel{=_{\mathrm{df}}}\;\sim(\sim\alpha\lor\sim\beta)
\]

\section{定理と証明}

\begin{theorem}{PM-Th.1（推移律）}{th1}
\((q\supset r)\supset\{(p\supset q)\supset(p\supset r)\}\)
\end{theorem}
\begin{tproof}{PM-Th.1}{pf1}
公理\(\mathrm{A4}\)に一様代入を施し，\(\supset\)の定義で読み替えるだけで得られる。
\begin{align*}
  1.\quad & (p\supset q)\supset\{(r\lor p)\supset(r\lor q)\} && [\mathrm{A4}]\\
  2.\quad & (q\supset r)\supset\{(\sim p\lor q)\supset(\sim p\lor r)\}
    && [1,\mathrm{PDR1}: p\mapsto q,\ q\mapsto r,\ r\mapsto\sim p]\\
  3.\quad & (q\supset r)\supset\{(p\supset q)\supset(p\supset r)\}
    && [2,\ \supset\text{の定義}]
\end{align*}
\end{tproof}

\begin{remark}{派生規則 PDR3（連鎖律）}{pdr3}
\(\mathrm{PM}\vdash\alpha\supset\beta\) かつ \(\mathrm{PM}\vdash\beta\supset\gamma\) ならば
\(\mathrm{PM}\vdash\alpha\supset\gamma\)。

PM-Th.1を一様代入した式 \((\beta\supset\gamma)\supset\{(\alpha\supset\beta)\supset(\alpha\supset\gamma)\}\)
に対し，\(\beta\supset\gamma\)でPDR2を適用して\((\alpha\supset\beta)\supset(\alpha\supset\gamma)\)を得て，
さらに\(\alpha\supset\beta\)でPDR2を適用すれば\(\alpha\supset\gamma\)が得られる（PDR2を2回）。
\end{remark}

\begin{theorem}{PM-Th.2（反射律）}{th2}
\(p\supset p\)
\end{theorem}
\begin{tproof}{PM-Th.2}{pf2}
\begin{align*}
  1.\quad & (q\supset r)\supset\{(p\supset q)\supset(p\supset r)\} && [\mathrm{PM\text{-}Th.1}]\\
  2.\quad & (p\lor q\supset p)\supset\{(p\supset p\lor q)\supset(p\supset p)\}
    && [1,\mathrm{PDR1}]\\
  3.\quad & p\lor p\supset p && [\mathrm{A1}]\\
  4.\quad & (p\supset p\lor p)\supset(p\supset p) && [2,3,\mathrm{PDR2}]\\
  5.\quad & p\supset p\lor p && [\mathrm{A2}]\\
  6.\quad & p\supset p && [4,5,\mathrm{PDR2}]
\end{align*}
\end{tproof}

\begin{theorem}{PM-Th.3（排中律）}{th3}
\(p\lor\sim p\)
\end{theorem}
\begin{tproof}{PM-Th.3}{pf3}
\(p\supset p\)を\(\supset\)の定義で読み替えると\(\sim p\lor p\)である。これに交換律A3を適用する。
\begin{align*}
  1.\quad & p\supset p\quad(\text{定義により }\sim p\lor p) && [\mathrm{PM\text{-}Th.2}]\\
  2.\quad & p\lor q\supset q\lor p && [\mathrm{A3}]\\
  3.\quad & \sim p\lor p\supset p\lor\sim p && [2,\mathrm{PDR1}: p\mapsto\sim p,\ q\mapsto p]\\
  4.\quad & p\lor\sim p && [1,3,\mathrm{PDR2}]
\end{align*}
\end{tproof}

\begin{theorem}{PM-Th.4a（二重否定の導入）}{th4a}
\(p\supset\sim\sim p\)
\end{theorem}
\begin{tproof}{PM-Th.4a}{pf4a}
排中律PM-Th.3に\(p\mapsto\sim p\)の一様代入を施す。
\begin{align*}
  1.\quad & p\lor\sim p && [\mathrm{PM\text{-}Th.3}]\\
  2.\quad & \sim p\lor\sim\sim p && [1,\mathrm{PDR1}: p\mapsto\sim p]\\
  3.\quad & p\supset\sim\sim p && [2,\ \supset\text{の定義}]
\end{align*}
\end{tproof}

\begin{remark}{派生規則 DR4（選言肢の付加）}{dr4}
\(\mathrm{PM}\vdash\alpha\supset\beta\) ならば \(\mathrm{PM}\vdash(\gamma\lor\alpha)\supset(\gamma\lor\beta)\)。

公理A4に\(p\mapsto\alpha,\ q\mapsto\beta,\ r\mapsto\gamma\)の一様代入を施すと
\((\alpha\supset\beta)\supset\{(\gamma\lor\alpha)\supset(\gamma\lor\beta)\}\)。
仮定\(\alpha\supset\beta\)にPDR2を適用すれば\((\gamma\lor\alpha)\supset(\gamma\lor\beta)\)が得られる。
\end{remark}

\begin{theorem}{PM-Th.4b（二重否定の除去）}{th4b}
\(\sim\sim p\supset p\)
\end{theorem}
\begin{tproof}{PM-Th.4b}{pf4b}
DR4とPDR3を使って，二重否定の導入（Th4a）から除去を導く。
\begin{align*}
  1.\quad & p\supset\sim\sim p && [\mathrm{PM\text{-}Th.4a}]\\
  2.\quad & \sim p\supset\sim\sim\sim p && [1,\mathrm{PDR1}: p\mapsto\sim p]\\
  3.\quad & (\sim p\supset\sim\sim\sim p)\supset\{(r\lor\sim p)\supset(r\lor\sim\sim\sim p)\}
    && [\mathrm{A4},\mathrm{PDR1}: p\mapsto\sim p,\ q\mapsto\sim\sim\sim p]\\
  4.\quad & (r\lor\sim p)\supset(r\lor\sim\sim\sim p) && [2,3,\mathrm{PDR2}]\\
  5.\quad & (p\lor\sim p)\supset(p\lor\sim\sim\sim p) && [4,\mathrm{PDR1}: r\mapsto p]\\
  6.\quad & p\lor\sim p && [\mathrm{PM\text{-}Th.3}]\\
  7.\quad & p\lor\sim\sim\sim p && [5,6,\mathrm{PDR2}]\\
  8.\quad & p\lor q\supset q\lor p && [\mathrm{A3}]\\
  9.\quad & (p\lor\sim\sim\sim p)\supset(\sim\sim\sim p\lor p) && [8,\mathrm{PDR1}: q\mapsto\sim\sim\sim p]\\
  10.\quad & \sim\sim\sim p\lor p && [7,9,\mathrm{PDR2}]\\
  11.\quad & \sim\sim p\supset p && [10,\ \supset\text{の定義：}\sim(\sim\sim p)\lor p=_{\mathrm{df}}\sim\sim p\supset p]
\end{align*}
（手順4〜7がDR4の実体，8〜10が交換律の適用である。）
\end{tproof}

\begin{theorem}{PM-Th.5a（移出律・対偶）}{th5a}
\((p\supset q)\supset(\sim q\supset\sim p)\)
\end{theorem}
\begin{tproof}{PM-Th.5a}{pf5a}
\begin{align*}
  1.\quad & p\supset\sim\sim p && [\mathrm{PM\text{-}Th.4a}]\\
  2.\quad & q\supset\sim\sim q && [1,\mathrm{PDR1}: p\mapsto q]\\
  3.\quad & (p\supset q)\supset\{(r\lor p)\supset(r\lor q)\} && [\mathrm{A4}]\\
  4.\quad & (q\supset\sim\sim q)\supset\{(\sim p\lor q)\supset(\sim p\lor\sim\sim q)\}
    && [3,\mathrm{PDR1}: p\mapsto q,\ q\mapsto\sim\sim q,\ r\mapsto\sim p]\\
  5.\quad & (\sim p\lor q)\supset(\sim p\lor\sim\sim q) && [2,4,\mathrm{PDR2}]\\
  6.\quad & p\lor q\supset q\lor p && [\mathrm{A3}]\\
  7.\quad & (\sim p\lor\sim\sim q)\supset(\sim\sim q\lor\sim p)
    && [6,\mathrm{PDR1}: p\mapsto\sim p,\ q\mapsto\sim\sim q]\\
  8.\quad & (\sim p\lor q)\supset(\sim\sim q\lor\sim p) && [5,7,\mathrm{PDR3}]\\
  9.\quad & (p\supset q)\supset(\sim q\supset\sim p)
    && [8,\ \supset\text{の定義（左辺・右辺とも）}]
\end{align*}
（手順1〜5がTh4aをA4で「\(\lor\)の中に押し込む」操作，6〜8が交換律，9で両辺を
\(\supset\)の定義に戻している。）
\end{tproof}

\begin{theorem}{PM-Th.6a（ド・モルガンの一方向）}{th6a}
\(\sim(p\land q)\supset(\sim p\lor\sim q)\)
\end{theorem}
\begin{tproof}{PM-Th.6a}{pf6a}
\(\land\)の定義\(p\land q=_{\mathrm{df}}\sim(\sim p\lor\sim q)\)より，
\(\sim(p\land q)\)は\(\sim\sim(\sim p\lor\sim q)\)と同じ式である。
これはPM-Th.4bの\(p\)に\((\sim p\lor\sim q)\)を代入した式そのものである。
\begin{align*}
  1.\quad & \sim\sim p\supset p && [\mathrm{PM\text{-}Th.4b}]\\
  2.\quad & \sim\sim(\sim p\lor\sim q)\supset(\sim p\lor\sim q)
    && [1,\mathrm{PDR1}: p\mapsto(\sim p\lor\sim q)]\\
  3.\quad & \sim(p\land q)\supset(\sim p\lor\sim q)
    && [2,\ \land\text{の定義：}p\land q=_{\mathrm{df}}\sim(\sim p\lor\sim q)]
\end{align*}
\end{tproof}

\section{まとめ}

以上で，公理A1--A4と原始規則PDR1・PDR2のみから，派生規則PDR3・DR4を経由して，
PM-Th.1〜PM-Th.6aまでを構文論的にすべて証明した。依存関係は次の通りである。
\[
  \mathrm{A1\text{--}A4}
  \;\to\;
  \mathrm{Th.1}
  \;\to\;
  \begin{cases}
    \mathrm{PDR3} \to \mathrm{Th.5a\ の一部}\\
    \mathrm{Th.2} \to \mathrm{Th.3} \to \mathrm{Th.4a}
    \to
    \begin{cases}
      \mathrm{DR4} \to \mathrm{Th.4b} \to \mathrm{Th.6a}\\
      \mathrm{Th.5a}
    \end{cases}
  \end{cases}
\]

\end{document}
